O sistema fuzzy Mamdani desenvolvido mostrou-se robusto, flexível e interpretável, sendo capaz de fornecer recomendações qualitativas de ventilação adequadas ao contexto de conforto térmico e eficiência energética. O método foi fundamentado teoricamente, detalhado em suas etapas e validado por meio de resultados gráficos e quantitativos, demonstrando sua eficácia para tomada de decisão sob incerteza em ambientes reais.

Por exemplo, no cenário 1 (temperatura de 28°C, umidade de 80\% e 15 pessoas), o sistema recomendou ventilação \textbf{FORTE}, com valor centroide de 62,44, refletindo corretamente a necessidade de maior renovação do ar em ambientes quentes, úmidos e cheios.

A abordagem fuzzy permite fácil adaptação a diferentes contextos e requisitos, além de lidar bem com incertezas e variações das condições ambientais. O estudo reforça a importância da escolha adequada das funções de pertinência, operadores e técnicas de defuzzificação para o desempenho do sistema.

% Como trabalhos futuros, sugere-se:
% \begin{itemize}
%     \item Aplicar o sistema fuzzy em dados reais de ambientes monitorados;
%     \item Explorar sistemas fuzzy com múltiplas entradas e saídas~\cite{ross2010fuzzy};
%     \item Investigar técnicas avançadas de otimização dos parâmetros fuzzy~\cite{haykin2008neural};
%     \item Integrar o sistema a plataformas de automação residencial ou industrial.
% \end{itemize}